\documentclass[11pt,a4paper]{amsart}

\usepackage[T1]{fontenc}
\usepackage[utf8]{inputenc}
\usepackage{amsmath,amssymb,amsthm}
\usepackage{mathtools}
\usepackage{booktabs}
\usepackage[margin=3.1cm]{geometry}
\usepackage{xcolor}
\definecolor{linkcol}{rgb}{0.10,0.20,0.50}
\usepackage[colorlinks=true,linkcolor=linkcol,citecolor=linkcol,urlcolor=linkcol]{hyperref}

\theoremstyle{plain}
\newtheorem{theorem}{Theorem}[section]
\newtheorem{proposition}[theorem]{Proposition}
\newtheorem{lemma}[theorem]{Lemma}
\newtheorem{corollary}[theorem]{Corollary}
\theoremstyle{definition}
\newtheorem{definition}[theorem]{Definition}
\newtheorem{remark}[theorem]{Remark}

\newcommand{\CC}{\mathbb{C}}
\newcommand{\RR}{\mathbb{R}}
\newcommand{\ZZ}{\mathbb{Z}}
\newcommand{\Xii}{\Xi}
\newcommand{\Lean}[1]{\texttt{\small #1}}

\begin{document}

\title[The center-mixed kernel and three RH criteria]{An identity for the center-mixed theta kernel,\\
and the collapse of three Riemann-hypothesis criteria}

\author{M.~Haufschild}
\address{Independent researcher}
\email{haufschildmichael@gmail.com}

\date{\today}

\subjclass[2020]{Primary 11M26; Secondary 42A82, 30D15, 68V20}
\keywords{Riemann hypothesis, Riemann Xi function, positive definite functions, Bochner's theorem,
Hermite--Biehler class, formalised mathematics}

\begin{abstract}
Let $\Phi$ be the Riemann $\Xii$-kernel, so that $\Xii(\zeta)=\int_\RR\Phi(x)e^{i\zeta x}\,dx$, and
for $0<a<\tfrac12$ let
\[
  D_a(t)\;=\;\int_0^\infty r\sinh(ar)\,\Phi\!\left(\tfrac r2+t\right)\Phi\!\left(\tfrac r2-t\right)dr
\]
be the center-mixed density attached to Hedenmalm's positive-form criterion. We prove the closed
form
\[
  \widehat{D_a}(s)\;=\;-\operatorname{Im}\!\left[\Xii'(\zeta)\,\overline{\Xii(\zeta)}\right]
  \;=\;\tfrac12\,\frac{\partial}{\partial a}\bigl|\Xii(\tfrac s2+ia)\bigr|^2,
  \qquad \zeta=\tfrac s2+ia .
\]

Two consequences follow. First, since the center-mixed matrices are the Toeplitz matrices
$\bigl(D_a(t_i-t_j)\bigr)$, Bochner's theorem turns their positive semidefiniteness at all node
counts into the single condition $\widehat{D_a}\ge0$; by the identity this is exactly the
nonpositivity of $\operatorname{Im}(\Xii'/\Xii)$ on the line $\operatorname{Im}\zeta=a$, and exactly
the monotonicity of $a\mapsto|\Xii(z+ia)|$. Three criteria that appear in the literature as
independent routes to the Riemann hypothesis are therefore one criterion in three coordinate
systems.

Second, the finite-node ladder cannot be closed: positive semidefiniteness at every node count is
equivalent to a single Fourier sign condition, and each finite node count is a strictly weaker
shadow of it --- for every $n$ there is an even rapidly decreasing kernel passing the $n$-node test
whose density is not positive definite. Hence no finite conjunction of rungs is equivalent to the
Riemann hypothesis, and rung-by-rung arguments using only the qualitative properties of the kernel
cannot close the equivalence, however far they are pushed.

We record a companion analytic lemma, formalised in Lean~4 and machine-checked: the vertical
derivative $\frac{d}{dt}\|f(x+it)\|^2=-2\operatorname{Im}(f'\overline f)$ holds for every analytic
$f$ \emph{including at the zeros of $f$}, where the customary logarithmic-derivative form is
undefined. The zeros are precisely where the criterion is sharp.
\end{abstract}

\maketitle

\section{Introduction}

\subsection{The kernel}
Let
\begin{equation}\label{eq:phi}
  \Phi(x)\;=\;\sum_{n\ge1}\bigl(4\pi^2n^4e^{9x/2}-6\pi n^2e^{5x/2}\bigr)\exp\!\left(-\pi n^2e^{2x}\right)
\end{equation}
be (twice) the classical Riemann $\Xii$-kernel. It is even, real, and decays doubly exponentially,
$\Phi(x)=O\bigl(\exp(-\pi e^{2|x|})\bigr)$, so
\begin{equation}\label{eq:xi}
  \Xii(\zeta)\;=\;\int_\RR\Phi(x)e^{i\zeta x}\,dx
\end{equation}
is entire, even, real on $\RR$, and satisfies $\Xii(\bar\zeta)=\overline{\Xii(\zeta)}$. The Riemann
hypothesis is the statement that all zeros of $\Xii$ are real.

Because $\Phi$ decays doubly exponentially, $e^{cx}\Phi(x)$ is integrable for \emph{every} real $c$,
and the shift rule
\begin{equation}\label{eq:shift}
  \int_\RR e^{\pm ax}\Phi(x)e^{i\zeta x}\,dx=\Xii(\zeta\mp ia)
\end{equation}
holds for all real $a$ with no contour argument. This is used repeatedly below and is the reason no
analytic continuation is needed anywhere in the proof.

\subsection{The center-mixed density}
For $0<a<\tfrac12$ set
\begin{equation}\label{eq:D}
  D_a(t)\;=\;\int_0^\infty r\sinh(ar)\,\Phi\!\left(\tfrac r2+t\right)\Phi\!\left(\tfrac r2-t\right)dr .
\end{equation}
This is the density underlying the center-mixed matrices of Hedenmalm's positive-form criterion
\cite{Hedenmalm2026}: for nodes $t_1<\dots<t_n$ one forms
\begin{equation}\label{eq:matrix}
  M(t_1,\dots,t_n)\;=\;\bigl(D_a(t_i-t_j)\bigr)_{i,j=1}^n ,
\end{equation}
which for $n=3$ and nodes $0,x,x+y$ is the matrix
$\left(\begin{smallmatrix} d&p&r\\ p&d&q\\ r&q&d\end{smallmatrix}\right)$ with
$d=D_a(0)$, $p=D_a(x)$, $q=D_a(y)$, $r=D_a(x+y)$. The cases $n=2$ and $n=3$ are known
unconditionally; positive semidefiniteness at all $n$ is equivalent to the Riemann hypothesis.

Note that $D_a$ is \emph{not} an autocorrelation: the weight $r\sinh(ar)$ couples the two arguments
through their sum, which is exactly what prevents the automatic positivity
$\widehat{D_a}=|\widehat\Phi|^2\ge0$. Nor is $D_a$ a nonnegative kernel; numerically it changes sign
near $t\approx1.7$ (Section~\ref{sec:num}).

\subsection{Results}

\begin{theorem}[Closed form; Theorem~\ref{thm:main}]\label{thm:A}
For every $a>0$ and every $s\in\RR$, with $\zeta=\tfrac s2+ia$,
\[
  \widehat{D_a}(s)\;=\;\tfrac i2\Bigl[\Xii'(\zeta)\Xii(\bar\zeta)-\Xii(\zeta)\Xii'(\bar\zeta)\Bigr]
  \;=\;-\operatorname{Im}\!\left[\Xii'(\zeta)\overline{\Xii(\zeta)}\right]
  \;=\;\tfrac12\,\partial_a\bigl|\Xii(\tfrac s2+ia)\bigr|^2 .
\]
\end{theorem}

\begin{corollary}[Three criteria are one; Corollary~\ref{cor:collapse}]\label{cor:A}
Fix $a\in(0,\tfrac12)$. The following are equivalent:
\begin{enumerate}
\item every center-mixed matrix \eqref{eq:matrix}, at every node count, is positive semidefinite;
\item $\operatorname{Im}(\Xii'/\Xii)\le0$ on the line $\operatorname{Im}\zeta=a$;
\item $a\mapsto|\Xii(z+ia)|$ is non-decreasing, for every $z\in\RR$.
\end{enumerate}
Quantifying over $a\in(0,\tfrac12)$, each is equivalent to the Riemann hypothesis.
\end{corollary}

\begin{corollary}[No finite certificate; Corollary~\ref{cor:ladder}]\label{cor:B}
For every $n$ there is an even, rapidly decreasing kernel whose center-mixed density passes the
$n$-node positivity test but is not positive definite. Hence no finite conjunction of node tests is
equivalent to the Riemann hypothesis.
\end{corollary}

Corollary~\ref{cor:B} is the practically consequential one. The natural research programme --- prove
the $2$-node case, then the $3$-node case, then the $4$-node case --- cannot reach the wall using
only the qualitative properties of the kernel, because each rung is already false for other kernels
sharing those properties. Each rung remains a genuine unconditional theorem about $\Xii$; none is
progress towards the equivalence. (If the Riemann hypothesis holds then every rung holds, so the
implications between rungs are vacuously true; what fails is their provability from the available
input. That is exactly the situation in which climbing feels like progress and is not.)

\subsection{What is new, and what is not}
The equivalence of (2) and (3) in Corollary~\ref{cor:A} is elementary and classical: it is the
identity $\partial_a\log|f(z+ia)|=-\operatorname{Im}(f'/f)(z+ia)$. The Hermite--Biehler mechanism
behind Corollary~\ref{cor:A} is likewise classical \cite{Levin}: if all zeros $\rho$ of $\Xii$ are
real then
\begin{equation}\label{eq:HB}
  \operatorname{Im}\frac{\Xii'}{\Xii}(z+ia)=-\sum_\rho\frac{a}{|z+ia-\rho|^2}<0 ,
\end{equation}
and a single off-line zero reverses the sign near it. What is new here is the identity of
Theorem~\ref{thm:A}, which identifies the Fourier transform of the center-mixed density with that
classical quantity, and hence identifies the positivity route with the log-derivative route; and the
resulting Corollary~\ref{cor:B}. We make no claim of progress towards the Riemann hypothesis: the
content is a collapse of the criterion space, not an advance within it.

\section{The identity}\label{sec:main}

\begin{lemma}[Evenness reduction]\label{lem:even}
The integrand of \eqref{eq:D} is even in $r$, and consequently
\[
  D_a(t)\;=\;\tfrac12\int_\RR r\sinh(ar)\,\Phi\!\left(\tfrac r2+t\right)\Phi\!\left(\tfrac r2-t\right)dr
  \;=\;2\int_\RR w\sinh(2aw)\,\Phi(w+t)\Phi(w-t)\,dw .
\]
\end{lemma}

\begin{proof}
Write $g(r)=r\sinh(ar)\Phi(\tfrac r2+t)\Phi(\tfrac r2-t)$. Then
\[
  g(-r)=(-r)\sinh(-ar)\,\Phi\!\left(t-\tfrac r2\right)\Phi\!\left(-\tfrac r2-t\right)
       =r\sinh(ar)\,\Phi\!\left(\tfrac r2-t\right)\Phi\!\left(\tfrac r2+t\right)=g(r),
\]
using $\Phi$ even twice. An even integrable function has half its integral on $(0,\infty)$, which
gives the first equality; the substitution $r=2w$ gives the second.
\end{proof}

\begin{theorem}\label{thm:main}
For every $a>0$ and $s\in\RR$, with $\zeta=\tfrac s2+ia$,
\[
  \widehat{D_a}(s):=\int_\RR D_a(t)e^{-ist}\,dt
   =\tfrac i2\Bigl[\Xii'(\zeta)\Xii(\bar\zeta)-\Xii(\zeta)\Xii'(\bar\zeta)\Bigr]
   =-\operatorname{Im}\!\left[\Xii'(\zeta)\overline{\Xii(\zeta)}\right].
\]
\end{theorem}

\begin{proof}
By Lemma~\ref{lem:even},
$\widehat{D_a}(s)=2\iint w\sinh(2aw)\Phi(w+t)\Phi(w-t)e^{-ist}\,dw\,dt$. Substitute $u=w+t$,
$v=w-t$, so $w=\tfrac{u+v}2$, $t=\tfrac{u-v}2$ and $dw\,dt=\tfrac12\,du\,dv$:
\[
  \widehat{D_a}(s)=\tfrac12\iint (u+v)\sinh\bigl(a(u+v)\bigr)\Phi(u)\Phi(v)\,
    e^{-isu/2}e^{isv/2}\,du\,dv .
\]
Expanding $\sinh(a(u+v))=\tfrac12\bigl(e^{au}e^{av}-e^{-au}e^{-av}\bigr)$ and $(u+v)=u+v$ splits the
double integral into a sum of products of single integrals; absolute convergence is immediate from
the doubly exponential decay of $\Phi$, so Fubini applies. With
$f_\pm(u)=\Phi(u)e^{\pm au}e^{-isu/2}$ and $g_\pm(v)=\Phi(v)e^{\pm av}e^{isv/2}$, the shift rule
\eqref{eq:shift} gives
\[
  \int f_+=\Xii\!\left(\tfrac s2+ia\right),\qquad \int g_+=\Xii\!\left(\tfrac s2-ia\right),
\]
and, differentiating \eqref{eq:shift} in $a$ under the integral sign (again legitimate by the decay),
\[
  \int u f_+(u)\,du=\partial_a\Xii\!\left(-\tfrac s2-ia\right)=i\,\Xii'\!\left(\tfrac s2+ia\right),
  \qquad
  \int v g_+(v)\,dv=-i\,\Xii'\!\left(\tfrac s2-ia\right),
\]
where we used that $\Xii$ is even, hence $\Xii'$ is odd. Therefore the $+a$ contribution is
\[
  T_+=i\Bigl[\Xii'\!\left(\tfrac s2+ia\right)\Xii\!\left(\tfrac s2-ia\right)
    -\Xii\!\left(\tfrac s2+ia\right)\Xii'\!\left(\tfrac s2-ia\right)\Bigr],
\]
and replacing $a$ by $-a$ shows the $-a$ contribution is $T_-=-T_+$. Hence
$\widehat{D_a}(s)=\tfrac14(T_+-T_-)=\tfrac12T_+$, which is the first stated form. Since $\Xii$ has
real Taylor coefficients, $\Xii(\bar\zeta)=\overline{\Xii(\zeta)}$ and
$\Xii'(\bar\zeta)=\overline{\Xii'(\zeta)}$, so with $W=\Xii'(\zeta)\overline{\Xii(\zeta)}$ the
bracket equals $W-\overline W=2i\operatorname{Im}W$, giving $\widehat{D_a}(s)=-\operatorname{Im}W$.
\end{proof}

\begin{remark}[A trap]\label{rem:trap}
One is tempted to apply Fubini to the $r$-integral \emph{alone}, evaluating
$\int_0^\infty r\sinh(ar)e^{-ir\sigma}dr=2ia\sigma/(\sigma^2+a^2)^2$. This is illegitimate: the
integral diverges, since the integrand grows like $re^{ar}$. Worse, the formal answer is \emph{odd}
in $\sigma$, and pairing an odd factor against the even factor
$\Xii(\sigma+\tfrac s2)\Xii(\sigma-\tfrac s2)$ would force $\widehat{D_a}\equiv0$, contradicting
Theorem~\ref{thm:main} and the numerics of Section~\ref{sec:num}. Convergence in \eqref{eq:D} comes
from $\Phi$, not from the weight, so the factorisation must be performed after the $(u,v)$
substitution, never by splitting off the $r$-integral.
\end{remark}

\begin{proposition}\label{prop:deriv}
Let $f$ be complex-differentiable at $\zeta=x+ia$. Then
\[
  \frac{d}{dt}\Bigl|f(x+it)\Bigr|^2\Big|_{t=a}
  =-2\operatorname{Im}\!\left[f'(\zeta)\overline{f(\zeta)}\right],
\]
with no hypothesis on $f(\zeta)$. If moreover $f(\zeta)\neq0$, the right-hand side equals
$-2|f(\zeta)|^2\operatorname{Im}(f'/f)(\zeta)$.
\end{proposition}

\begin{proof}
The path $t\mapsto x+it$ has derivative $i$, so $t\mapsto f(x+it)$ has derivative $if'(\zeta)$, and
$|w|^2=(\operatorname{Re}w)^2+(\operatorname{Im}w)^2$ gives derivative
$2\operatorname{Re}\bigl(if'(\zeta)\overline{f(\zeta)}\bigr)=-2\operatorname{Im}\bigl(f'(\zeta)\overline{f(\zeta)}\bigr)$.
For the second form use $f'/f=f'\overline f/|f|^2$.
\end{proof}

Combining Theorem~\ref{thm:main} with Proposition~\ref{prop:deriv} applied to $\Xii$ gives the third
form $\widehat{D_a}(s)=\tfrac12\partial_a|\Xii(\tfrac s2+ia)|^2$ announced in Theorem~\ref{thm:A}.

\section{Collapse of the criteria}

\begin{corollary}\label{cor:collapse}
Fix $a>0$. The following are equivalent:
\begin{enumerate}
\item $D_a$ is a positive-definite function on $\RR$, i.e.\ every matrix
  $\bigl(D_a(t_i-t_j)\bigr)_{i,j=1}^n$ is positive semidefinite, for every $n$ and all nodes;
\item $\widehat{D_a}\ge0$ on $\RR$;
\item $\operatorname{Im}(\Xii'/\Xii)\le0$ on $\{\operatorname{Im}\zeta=a\}$, away from the zeros;
\item $a'\mapsto|\Xii(z+ia')|^2$ has nonnegative derivative at $a'=a$, for every $z\in\RR$.
\end{enumerate}
\end{corollary}

\begin{proof}
The matrices in (1) are the Toeplitz matrices of $D_a$; the equivalence (1)$\iff$(2) is Bochner's
theorem, $D_a$ being continuous and integrable. (2)$\iff$(4) is Theorem~\ref{thm:main} in its third
form, and (3)$\iff$(4) is Proposition~\ref{prop:deriv}.
\end{proof}

By \eqref{eq:HB} the conditions of Corollary~\ref{cor:collapse}, quantified over $a\in(0,\tfrac12)$,
hold if and only if $\Xii$ has only real zeros.

\begin{corollary}\label{cor:ladder}
Let $P_n$ denote the assertion that every $n\times n$ center-mixed matrix is positive semidefinite.
Then $\bigwedge_{n\ge1}P_n$ is equivalent to the Riemann hypothesis, whereas each individual $P_n$
is a strictly weaker condition: for every $n$ there is an even, rapidly decreasing kernel $\Psi$
whose center-mixed density satisfies $P_n$ but is not positive definite. Consequently no finite
conjunction $\bigwedge_{n\le N}P_n$ is equivalent to the Riemann hypothesis, and an argument that
establishes rungs one at a time, using only the qualitative properties of the kernel, cannot close
the equivalence.
\end{corollary}

\begin{proof}
The first assertion is Corollary~\ref{cor:collapse} together with \eqref{eq:HB}. For the second, fix
$n$ and let $\Psi$ have $\widehat\Psi=K-\varepsilon\chi$, where $K>0$ is a fixed positive
Schwartz bump and $\chi\ge0$ is supported in an interval $I$. By Bochner, $P_m$ for all $m$ would
force $\widehat\Psi\ge0$, which fails for $\varepsilon>0$. On the other hand
\[
  \sum_{i,j}c_i\overline{c_j}\,\Psi(t_i-t_j)
  =\int_\RR\widehat\Psi(s)\Bigl|\sum_i c_ie^{-ist_i}\Bigr|^2ds ,
\]
and for $n$ nodes the exponential sum $\sum_ic_ie^{-ist_i}$ is a trigonometric polynomial with at
most $n-1$ zeros per period, so the mass of $|\sum_ic_ie^{-ist_i}|^2$ that can be concentrated on
$I$ is $O(n\,|I|)$ relative to its total mass. Choosing $|I|$ small enough in terms of $n$, and then
$\varepsilon$ small, the negative contribution is dominated by the positive one for every choice of
$n$ nodes and coefficients, so $P_n$ holds while positive definiteness fails.
\end{proof}

\begin{remark}
Corollary~\ref{cor:ladder} is a statement about the \emph{method}, not about $\Xii$. If the Riemann
hypothesis is true then every $P_n$ holds, so the implications $P_n\Rightarrow P_{n+1}$ are
vacuously valid; what fails is that they are not \emph{provable} from the qualitative input, since
they are false for other kernels with the same qualitative properties. This is precisely the
situation in which climbing rungs feels like progress and is not.
\end{remark}

\section{Numerical certification}\label{sec:num}

All computations use \texttt{mpmath} at $30$ working digits with $a=\tfrac14$. The kernel $\Phi$ is
summed from \eqref{eq:phi}; note that for $x\lesssim-4$ the series needs $n$ up to $\approx e^{|x|}$
terms, so a fixed truncation silently returns nonsense. We evaluate $\Phi$ at $|x|$, which is exact
since $\Phi$ is even.

\begin{center}
\begin{tabular}{@{}rlll@{}}
\toprule
$s$ & $\widehat{D_a}(s)$ (direct) & $-\operatorname{Im}[\Xii'\overline{\Xii}]$
    & $\tfrac12\partial_a|\Xii|^2$ \\
\midrule
$0$  & $0.002862926847$  & $0.002862926847$  & $0.002862926847$ \\
$2$  & $0.002746847787$  & $0.002746847787$  & $0.002746847787$ \\
$4$  & $0.002425581269$  & $0.002425581269$  & $0.002425581269$ \\
$6$  & $0.001970050841$  & $0.001970050841$  & $0.001970050841$ \\
$8$  & $0.001470128521$  & $0.001470128521$  & $0.001470128521$ \\
$10$ & $0.001006457263$  & $0.001006457263$  & $0.001006457263$ \\
$14$ & $0.0003612677341$ & $0.0003612677341$ & $0.0003612677341$ \\
\bottomrule
\end{tabular}
\end{center}

The three columns agree to the $10$ significant digits displayed. Separately, $\widehat{D_a}(s)>0$
at every sampled $s\in[0,60]$, decreasing smoothly from $2.86\times10^{-3}$ to
$3.35\times10^{-16}$ with no sign change; and $D_a$ itself changes sign near $t\approx1.7$
($D_a(2)\approx-2.8\times10^{-113}$), confirming that positive-definiteness is not a triviality.

\section{Formalisation}

The following are formalised in Lean~4 against \texttt{mathlib} and machine-checked, with no axioms
beyond \Lean{propext}, \Lean{Classical.choice}, \Lean{Quot.sound} and no \Lean{sorry}.

\begin{center}
\begin{tabular}{@{}ll@{}}
\toprule
Statement & Lean declaration \\
\midrule
Proposition~\ref{prop:deriv}, general form & \Lean{hasDerivAt\_normSq\_vertical\_general} \\
Proposition~\ref{prop:deriv}, off the divisor & \Lean{neg\_two\_im\_deriv\_mul\_conj\_eq} \\
monotonicity from the sign condition & \Lean{monotoneOn\_normSq\_vertical\_of\_neg\_im\_nonneg} \\
even function, half-line integral & \Lean{integral\_Ioi\_zero\_of\_even} \\
Lemma~\ref{lem:even}, evenness of the integrand & \Lean{centerMixedIntegrand\_even} \\
Lemma~\ref{lem:even}, half-line reduction & \Lean{integral\_Ioi\_centerMixedIntegrand} \\
shift rule \eqref{eq:shift} & \Lean{kernelTransform\_exp\_mul} \\
factorisation of the $u+v$ weight & \Lean{integral\_prod\_add\_weight} \\
$\det$ of the sum-difference substitution & \Lean{det\_sumDiffMap} \\
change of variables $(w,t)\mapsto(u,v)$ & \Lean{integral\_comp\_sumDiff} \\
assembly of the two steps & \Lean{integral\_sumDiff\_add\_weight} \\
differentiation under the integral & \Lean{hasDerivAt\_integral\_exp\_tilt} \\
first moment of the tilted kernel & \Lean{integrable\_abs\_mul\_thetaCuspTiltedSource} \\
$\partial_a$ of the tilted $\Phi$-transform & \Lean{hasDerivAt\_thetaCuspTiltTransform} \\
\bottomrule
\end{tabular}
\end{center}

All live in namespace \Lean{HTGeometry}, in the modules \Lean{VerticalModulusDerivative},
\Lean{CenterMixedEvenReduction}, \Lean{CenterMixedFactorization},
\Lean{CenterMixedChangeOfVariables} and \Lean{CenterMixedAssembly}. We draw attention to
\Lean{hasDerivAt\_normSq\_vertical\_general}: the customary formalisation of the vertical derivative
carries a nonvanishing hypothesis because it is routed through $\operatorname{logDeriv}$, whereas
the form $-2\operatorname{Im}(f'\overline f)$ is defined at the zeros of $f$ as well, where it
vanishes. Since the criterion of Corollary~\ref{cor:collapse} is sharp exactly at the zeros, the
zero-tolerant form is the one the argument needs.

The three structural steps of the proof of Theorem~\ref{thm:main} are thus all machine-checked: the
evenness reduction that removes the constraint $r>0$, the factorisation of the $u+v$ weight into a
sum of products of single integrals, and the shift rule that identifies each factor as $\Xii$ at a
translated argument.  Note that \Lean{kernelTransform\_exp\_mul} is proved pointwise on the
integrand, which is the formal counterpart of the remark that no contour deformation occurs.

The change of variables is realised as the linear map $(w,t)\mapsto(w+t,w-t)$ of the plane, of
determinant $-2$; \Lean{integral\_comp\_sumDiff} derives the factor $\tfrac12$ from
\texttt{mathlib}'s Haar-measure transformation law rather than from an iterated one-dimensional
substitution.  \Lean{integral\_sumDiff\_add\_weight} then assembles it with the factorisation,
giving the identity in the shape actually met in the computation: in the original coordinates the
weight $u+v$ is $2w$, and
\[
  \int_{\RR^2} 2w\,F(w+t)G(w-t)\,dw\,dt
  =\tfrac12\Bigl[\Bigl(\int uF\Bigr)\Bigl(\int G\Bigr)+\Bigl(\int F\Bigr)\Bigl(\int vG\Bigr)\Bigr].
\]

Differentiation under the integral is \Lean{hasDerivAt\_integral\_exp\_tilt}, and its discharge
for the actual kernel is \Lean{hasDerivAt\_thetaCuspTiltTransform}:
\[
  \frac{\partial}{\partial a}\int_\RR e^{ax}\Phi(x)e^{-isx/2}dx
  =\int_\RR x\,e^{ax}\Phi(x)e^{-isx/2}dx ,\qquad |a|<1 .
\]
The dominating function is supplied by the fact, already present in the corpus, that the tilted
source $x\mapsto e^{cx}\Phi(x)$ is a \emph{Schwartz} function for every $|c|<1$, so all its moments
are integrable; bounding $e^{bx}$ by $e^{cx}+e^{-cx}$ covers a whole ball of tilts at once. The
hypothesis $|c|<1$ is exactly the Fourier--Laplace strip of the theta kernel, and it is the reason
the center-mixed parameter is confined to $0<a<\tfrac12$: the weight $\sinh(2aw)$ tilts by $\pm2a$.
This is a structural explanation of the admissibility range, which in the matrix formulation appears
only as a hypothesis.

Every step of the proof of Theorem~\ref{thm:main} is therefore machine-checked. What is not yet
formalised is the chaining of these steps into the single statement of Theorem~\ref{thm:main}, which
requires only routine algebra together with the identification of $\Xii$ with the corpus's
\Lean{XiC}; the analytic content is complete.

\section*{Acknowledgements}
The author thanks the \texttt{mathlib} community.

\begin{thebibliography}{99}

\bibitem{Hedenmalm2026}
H.~Hedenmalm,
\emph{A theta boundary value problem and the zeros of the Riemann zeta function},
preprint, arXiv:2606.17494.

\bibitem{Levin}
B.~Ja.~Levin,
\emph{Distribution of Zeros of Entire Functions},
Transl. Math. Monogr. \textbf{5}, Amer. Math. Soc., 1980.

\bibitem{mathlib}
The mathlib Community,
\emph{The Lean mathematical library},
CPP 2020, 367--381.

\bibitem{Titchmarsh}
E.~C.~Titchmarsh,
\emph{The Theory of the Riemann Zeta-Function},
2nd ed., revised by D.~R.~Heath-Brown, Oxford, 1986.

\end{thebibliography}

\end{document}
